\appendix

\section{Inventario de artefactos y reproducibilidad}

La Tabla~\ref{tab:artefactos} enumera los artefactos principales del proyecto y su función en la cadena de evidencia. Todos los archivos de datos se conservan sin modificaciones respecto de la ejecución original; las tablas y figuras del informe se generan a partir de ellos mediante scripts reproducibles.

\begin{table*}[t]
\centering
\scriptsize
\setlength{\tabcolsep}{4pt}
\caption{Inventario de artefactos y trazabilidad de la reproducibilidad}
\label{tab:artefactos}
\begin{tabularx}{\textwidth}{@{} >{\raggedright\arraybackslash}p{3.0cm} >{\raggedright\arraybackslash}X >{\raggedright\arraybackslash}p{5.5cm} @{}}
\toprule
\textbf{Artefacto} & \textbf{Ruta} & \textbf{Contenido} \\
\midrule
Selección ampliada & \texttt{data/repos\_candidatos\_ampliados.csv} & 915 candidatos deduplicados. \\
Preselección & \texttt{data/repos\_preseleccion\_ampliada.csv} & 120 repositorios, 40 por banda. \\
Exclusiones & \texttt{data/repos\_excluidos\_ampliada.csv} & 124 registros con motivo de exclusión. \\
Elegibilidad & \texttt{data/decision\_muestra\_ampliada.csv} & 24 incluidos y 96 excluidos. \\
Relevancia local & \texttt{data/evidencia\_relevancia.csv} & 23 CUMPLE y 1 NO CUMPLE. \\
Muestra final & \texttt{data/repos\_seleccionados.csv} & 23 repositorios con SHA fijado. \\
Hallazgos & \texttt{data/resultados/findings.csv} & 445 hallazgos con archivo y línea. \\
Métricas & \texttt{data/resultados/metricas.csv} & 15\,768 módulos y cinco métricas. \\
Errores de métricas & \texttt{data/resultados/metricas\_errors.csv} & 1 error de sintaxis documentado. \\
Resumen de escaneo & \texttt{data/resultados/scan\_summary.csv} & 15\,769 archivos, 0 errores de análisis. \\
Contrastes & \texttt{data/resultados/estadisticas.csv} & 3 primarios, 9 por tercil y 3 de sensibilidad. \\
Correlaciones & \texttt{data/resultados/spearman.csv} & 3 correlaciones de Spearman. \\
Decisiones & \texttt{data/resultados/decisiones\_hipotesis.csv} & D9 y D10 apoyadas en el piloto. \\
Adjudicación IA & \texttt{data/resultados/validacion\_investigador.csv} & 60 casos con procedencia asistida por IA. \\
Figuras & \texttt{scripts/08\_generar\_figuras\_informe.py} & 8 figuras reproducibles en PNG. \\
Tablas & \texttt{scripts/09\_generar\_tablas\_latex.py} & Tablas LaTeX generadas desde los CSV. \\
\bottomrule
\end{tabularx}
\end{table*}

\subsection{Secuencia de reproducción}

La secuencia completa se ejecuta desde la raíz del proyecto. Los scripts de recolección requieren credenciales de GitHub y no se ejecutaron de nuevo para este informe; los scripts de análisis, figuras y tablas pueden ejecutarse sobre los datos ya conservados. La secuencia es la siguiente:

\begin{verbatim}
# Detección y metricas (requiere los clones locales)
python scripts/03_ejecutar_detectores.py
python scripts/04_calcular_metricas.py

# Analisis estadistico
python scripts/05_analisis_estadistico.py

# Muestreo y adjudicacion asistida por IA
python scripts/06_muestra_validacion.py
python scripts/07_adjudicar_validacion_asistida.py

# Figuras y tablas del informe
python scripts/08_generar_figuras_informe.py
python scripts/09_generar_tablas_latex.py
\end{verbatim}

La compilación del documento se realizó con MiKTeX, mediante la secuencia clásica de cuatro pasos:

\begin{verbatim}
pdflatex main.tex
bibtex main
pdflatex main.tex
pdflatex main.tex
\end{verbatim}

\subsection{Nota sobre la procedencia de los datos}

Los archivos \texttt{findings.csv}, \texttt{metricas.csv}, \texttt{estadisticas.csv}, \texttt{spearman.csv}, \texttt{resumen\_descriptivo.csv}, \texttt{resumen\_repositorio.csv}, \texttt{decisiones\_hipotesis.csv} y \texttt{validacion\_investigador.csv} son las fuentes primarias de las cifras reportadas. Las tablas y figuras no introducen valores nuevos: agregan, ordenan o representan esos datos. Cualquier discrepancia entre una tabla y un archivo CSV debe resolverse a favor del CSV, y este informe no modificó los CSV para ajustarlos a la narrativa.

La revisión asistida por IA se conserva en \texttt{revision\_ia\_preliminar.csv} y en \texttt{validacion\_investigador.csv}. Esta última incluye las columnas \texttt{adjudication\_method} e \texttt{independent\_human\_verification}, que documentan que la clasificación fue adoptada como evaluación asistida por IA y que no hubo verificación humana independiente. Esa procedencia es parte integral de los datos y no debe eliminarse en futuras reutilizaciones.
